\scanpage{94}
une composante $D_\alpha=\overline{\{s_\alpha\}}$ de ce diviseur, n'est autre que sa \uline{multiplicité} en $D_\alpha$.

Ceci rappelé, on se demande donc de ce qui se passe pour les $s\in\bigcup D_\alpha$, \deleted{On} \inserted{et} \deleted{voit que} regardons ceux qui sont $\in D_\alpha$, $\alpha$ fixé, i.e. qui spécialisent un $s_\alpha$ fixé. On peut dire que la situation définit une extension $\mathcal E_{U_{s_\alpha}}$ de $\mathcal E_{k(s_\alpha)}$ par $\pi$ --- déduite, comme tantôt pour $S$ et $K$, d'une courbe \unsure{relative} $U_{D_\alpha}$ sur $D_\alpha$ --- ou $U_{S_\alpha}$ sur un ouvert régulier convenable de $D_\alpha$, si on préfère --- en passant au pt générique. Cette fois-ci, pour la fibre générique, on n'a plus des \deleted{\unsure{sections}} pts. rat. sur $k(s_\alpha)$ de $\unsure{U_{s_\alpha}|U_{S_\alpha}}$ --- on n'a qu'un pt à l'infini $i_\alpha\in X_{s_\alpha}\setminus U_{s_\alpha}$, rat sur $\unsure{k_\alpha}$. Mais on dispose \inserted{maintenant} d'un $L_i\subset\pi$ relatif à cet $i_\alpha$, justement, de façon ess. canonique,
\margin{93}
